
\documentclass{article}
\usepackage{colortbl}
\usepackage{makecell}
\usepackage{multirow}
\usepackage{supertabular}

\begin{document}

\newcounter{utterance}

\centering \large Interaction Transcript for game `cladder', experiment `full\_v1.5\_default', episode 3466 with qwen3.5{-}9b.
\vspace{24pt}

{ \footnotesize  \setcounter{utterance}{1}
\setlength{\tabcolsep}{0pt}
\begin{supertabular}{c@{$\;$}|p{.15\linewidth}@{}p{.15\linewidth}p{.15\linewidth}p{.15\linewidth}p{.15\linewidth}p{.15\linewidth}}
    \# & \multicolumn{2}{c}{Player} && \multicolumn{2}{c}{Game Master} \\
    \hline

    \theutterance \stepcounter{utterance}  
    & & & \multicolumn{4}{p{0.6\linewidth}}{
        \cellcolor[rgb]{0.9,0.9,0.9}{
            \makecell[{{p{\linewidth}}}]{
                \texttt{\tiny{[P1$\langle$GM]}}
                \texttt{You are an expert in causal inference. The following question is not a typical commonsense query, but rather a meticulously designed question created by a professor specializing in causal inference, intended to assess the students' mastery of the course content.} \\
\texttt{Imagine a self{-}contained, hypothetical world with only the following conditions, and without any unmentioned factors or causal relationships: Solar eclipse has a direct effect on Alice arriving to school. Traffic has a direct effect on Alice arriving to school. Method 1: We look at how solar eclipse correlates with Alice arriving to school case by case according to traffic. Method 2: We look directly at how solar eclipse correlates with Alice arriving to school in general. To understand how solar eclipse affects Alice arriving to school, is it more correct to use the Method 1 than Method 2?} \\
\texttt{Start your answer with 'Yes' or 'No', followed by additional reasoning or evidence to support your explanation.} \\
            }
        }
    }
     \\ \\

    \theutterance \stepcounter{utterance}  
    & \multicolumn{4}{p{0.6\linewidth}}{
        \cellcolor[rgb]{0.561,0.737,0.561}{
            \makecell[{{p{\linewidth}}}]{
                \texttt{\tiny{[P1$\rangle$GM]}}
                \texttt{No.} \\
\\ 
\texttt{In this hypothetical world, **Traffic** acts as a **confounder** because it has a direct effect on the outcome (Alice arriving to school) and is presumably associated with the treatment (Solar eclipse) in the context of the data collection or natural variation, or more critically, the question implies we are trying to isolate the specific causal effect of the eclipse while traffic is a varying factor influencing the outcome.} \\
\\ 
\texttt{However, the phrasing "according to traffic" in Method 1 suggests a **stratification** or **conditioning** approach. In causal inference, to estimate the causal effect of a variable \$X\$ (Solar eclipse) on \$Y\$ (Arriving to school) when there is a confounder \$Z\$ (Traffic), we must condition on \$Z\$ to block the backdoor path or to ensure we are comparing like{-}with{-}like.} \\
\\ 
\texttt{If we use **Method 2** (looking at the correlation in general), we are calculating the marginal correlation. If Traffic is correlated with the occurrence of a Solar Eclipse (e.g., perhaps eclipses happen on days with specific traffic patterns, or if the "Solar Eclipse" variable in the dataset is actually a proxy for a specific time of day that also dictates traffic), the marginal correlation will be biased by the confounding effect of Traffic. It would conflate the effect of the eclipse with the effect of traffic.} \\
\\ 
\texttt{**Method 1** attempts to control for the confounder by looking at the relationship "case by case according to traffic." By stratifying on Traffic (or conditioning on it), we isolate the effect of the Solar Eclipse within levels of constant Traffic. This aligns with the principle of **backdoor adjustment**: \$P(Y|do(X)) = \textbackslash{}sum\_z P(Y|X, Z=z)P(Z=z)\$. By examining the correlation within strata of \$Z\$ (Traffic), we remove the bias introduced by the confounder.} \\
\\ 
\texttt{Therefore, to understand the *causal* effect of the solar eclipse specifically, Method 1 is the correct approach because it controls for the confounding variable (Traffic), whereas Method 2 yields a biased estimate due to the omitted variable bias from Traffic.} \\
            }
        }
    }
    & & \\ \\

    \theutterance \stepcounter{utterance}  
    & & & \multicolumn{2}{p{0.3\linewidth}}{
        \cellcolor[rgb]{0.9,0.9,0.9}{
            \makecell[{{p{\linewidth}}}]{
                \texttt{\tiny{[GM$|$GM]}}
                \texttt{No} \\
            }
        }
    }
    & & \\ \\

    \theutterance \stepcounter{utterance}  
    & & & \multicolumn{2}{p{0.3\linewidth}}{
        \cellcolor[rgb]{0.9,0.9,0.9}{
            \makecell[{{p{\linewidth}}}]{
                \texttt{\tiny{[GM$|$GM]}}
                \texttt{no} \\
            }
        }
    }
    & & \\ \\

    \theutterance \stepcounter{utterance}  
    & & & \multicolumn{2}{p{0.3\linewidth}}{
        \cellcolor[rgb]{0.9,0.9,0.9}{
            \makecell[{{p{\linewidth}}}]{
                \texttt{\tiny{[GM$|$GM]}}
                \texttt{game\_result = WIN} \\
            }
        }
    }
    & & \\ \\

\end{supertabular}
}

\end{document}
